% v0.30 — ανεξάρτητο regression δοκίμιο callouts και αριθμημένων δηλώσεων.
% Ίδιο περιεχόμενο σε SCREEN και PRINT, χωρίς καμία άμεση αλλαγή χρώματος.
\chapter{Δοκιμή διδακτικών ενθέτων}
\BookSection{Ιστορικό ένθετο και διδακτικά κουτιά}
\begin{HistoryBox}[Ο Γαλιλαίος και η μέτρηση]
Η ιστορία της Φυσικής συνδέεται με την ανάπτυξη της πειραματικής
μεθόδου. Ο Γαλιλαίος μελέτησε συστηματικά την κίνηση και βοήθησε
να καθιερωθεί η μέτρηση ως βάση για την επιστημονική ερμηνεία.

Οι συσκευές, οι πειραματικές παρατηρήσεις και οι μαθηματικές σχέσεις
επέτρεψαν την ακριβέστερη περιγραφή των φυσικών φαινομένων.
\end{HistoryBox}
\begin{ThinkBox}
Ελέγχω τα δεδομένα και ξεκινώ από τον κατάλληλο φυσικό ορισμό.
Η μαθηματική έκφραση πρέπει να έχει σαφές φυσικό νόημα.
\end{ThinkBox}
\begin{AttentionBox}
Η μετατόπιση και το διάστημα δεν είναι γενικά το ίδιο μέγεθος.
\end{AttentionBox}
\BookSection{Μαθηματικές δηλώσεις χωρίς εικονίδια}
\begin{definitionbox}[Συνάρτηση πρώτου βαθμού]\label{def:trial}
Η συνάρτηση $f(x)=\alpha x+\beta$, με $\alpha\ne 0$, είναι
συνάρτηση πρώτου βαθμού.
\end{definitionbox}
\begin{theorembox}[Πρόσθεση ίσων ποσοτήτων]\label{thm:trial}
Αν $a=b$, τότε $a+c=b+c$ για κάθε πραγματικό $c$.
\end{theorembox}
Ο ορισμός \ref{def:trial} και το θεώρημα \ref{thm:trial}
διατηρούν τους ανεξάρτητους μετρητές τους.
